% ===================================================
% CHAPITRE 19 — PRÊTRES NÉS À OUROUX
% Pages imprimées 176-180 — vues 179-183
% ===================================================
\chapter{Prêtres nés à Ouroux}

% --- Page imprimée 176 — vue 179
NÉS  
A OUROUX

I-\underline{JOANNES DE MONTELATINO}

II-\underline{BENEDICTUS DE MONTELATINO} vivaient vers la fin du XIVème siècle ou au commencement du XVème. Ils étaient frères. Leurs noms nous ont été conservés dans le terrier de Vauxrenard où nous lisons au chapitre des reconnaissances de Guyot Ligier : « qui quondam fuerunt Joannes et Benedictus de Montelatino fratrum et presbyterum ».

Nous avons vu dans la liste des curés d'Ouroux un Joannes de Montalino vicaire perpétuel en 1449. Il est probable qu'il y a eu erreur de la part du copiste dans l'acte de fondation de la chapelle de Ste Croix, qui d'ailleurs fourmille de fautes d'orthographe et qu'il a lu Montalino au lieu de Montelatino.

Le terrier dont il est fait question plus haut est de 1450.

III-\underline{HUGUES ou HUGONIN TESTENOIRE} appartenait à la famille des Ligier alias Testenoire qui devait jouer un rôle si important dans notre paroisse. Il était fils de Guyot Ligier et de Flor de Laye, son épouse. Son frère Thomas continua la branche aînée d'Ouroux ; Michel acheta la seigneurie de Bacot à St-Christophe-la-Montagne et fonda la branche des seigneurs du même nom.

Le nom de Hugues Ligier, alias Testenoire, se retrouve deux fois dans le terrier d'Alloignet, signé de Fonte et rédigé en 1529. Il est désigné au f.59 verso, dans la reconnaissance de Benoît Chambenard et au f.61 verso, dans celle de Clément Peyraud, comme possédant un bois à Monlong, près des Loyttes. Il vivait vers 1529 et était prêtre.

% --- Page imprimée 177 — vue 180
IV-\underline{Anthoine de Lafont} était curé de Cenves en 1623. Il appartenait à la famille des De Fonte qui succéda à celle de De Bosco dans le notariat d'Ouroux et fut remplacé plus tard par celle des Testenoire.
    Claude de Lafont père de Anthoine, curé de Cenves, était notaire à Ouroux en 1582; il vivait encore en 1603. Il laissa quatre enfants : Timothée, Vincent qui fut notaire à Matour, Anthoinette qui fut mariée à Jehan Testenoire et Anthoine qui entra dans les Ordres. Son frère Timothée lui donna une maison située au bourg d'Ouroux sur la grandrue tendant du pont des Chassaignes à la croix du Rempart d'Ouroux, près du colombier. Il fut le grand oncle, probablement même le parrain de Anthoine De Lafont qui épousa Catherine Alacoque, cousine de la Bienheureuse Marguerite Marie.

V-\underline{Estienne de Lafont} neveu du précédent était fils de Timothée de Lafont et de Constance Bellor. Timothée était gendarme de la compagnie du duc de Nemours; il mourut en 1623. Estienne fut ordonné diacre en 1629; nous l'avons vu prêtre, vicaire perpétuel d'Ouroux en 1630

VI-\underline{Guillaume du Tel} ne nous est connu que par le terrier d'Alloignet Il est mentionné au folio 71, dans la reconnaissance de Claude Berthelon surnommé Monnet, comme possédant un pré et une verchère auprès du village actuel des Monnet. La reconnaissance est de 1528.

VII-\underline{Benoît Morin} vicaire perpétuel d'Ouroux en 1629 puis curé d'Emeringes où il eut pour successeur son neveu:

VIII-\underline{Jean Chrysostome Morin} dont nous avons parlé longuement au chapitre concernant les curés de la paroisse

IX-\underline{Claude de Grosbois} naquit le 28 Juin 1676. Il était fils de Jean Chrysostome de Grosbois, capitaine-enseigne au régiment de Certirane et de Marie de la Roche. Son parrain fut Claude de la Roche, escuier, sieur de Poncié; sa marraine, Benoîte Blanc, absente et remplacée par Françoise Testenoire. Il fut l'aîné de six enfants dont cinq filles.
    En 1703 il était au séminaire des R.R.P.P. Jésuites de Toulouse; il était docteur en droit civil et en droit canonique.
    Afin d'être ordonné prêtre, son père lui constitua comme titre patrimonial les biens fonds, à lui appartenant, situés sur la paroisse de Chiroubles, comprenant quatre vignerons et ayant un revenu de 600 li
% --- Page imprimée 178 — vue 181 (suite)
vres sauf à payer une aumône de 24 mesures de blé et de 3 asnées de bon vin, à tous les pauvres, le Jeudi-Saint. Il administra la paroisse d'Ouroux en 1732 après le départ de Mr Perrachon jusqu'à l'arrivée de Mr Guillin qui eut lieu l'année suivante.

Claude de Grosbois mourut le 28 Juin 1765 âgé de 90 ans et fut inhumé sur sa demande à l'entrée de l'église sous la galonnière.

X-\underline{Claude Marie Duligier Testenoire} fut le sixième enfant des époux Jean-Marie Duligier et Marguerite Reine Nonin. Il fut baptisé le 21 février 1749 et eut pour parrain Claude Nonin marchand de Varennes, paroisse de Mazilles ; sa marraine fut Claire Marie Babillon femme de Aymé Gabriel Verset notaire royal à Tramayes.

Entré dans les Ordres il fut d'abord curé de Cenves, puis de la Chapelle-sous-Dun vers 1788. Durant la tourmente révolutionnaire il se retira au fief des Micoux, sur la paroisse de Vauxrenard ; ce fief appartenait, à cette époque, à la famille des Testenoire. Il y mourut atteint de la folie des persécutions, à la suite des scènes d'horreur dont il avait été le témoin et peut-être la victime.

XI-\underline{Jean Trompier} appartenait à la famille des Trompier venue de Chiroubles à Ouroux, pour enseigner à lire aux enfants. Étudiant au commencement de la révolution, il fut forcé de rentrer dans sa famille. Le maire Louis Matray eut recours à lui pour rédiger les actes de l'état civil jusqu'au 23 Nivôse An II où Jean Trompier dut partir comme faisant partie de la réquisition des jeunes gens de 18 à 25 ans.

Néanmoins il fut ordonné prêtre après 1800, puis nommé curé de Monsols et enfin décédé curé de Cras dans le diocèse de Belley, le 18 Avril 1835.

Dès 1812 il avait rencontré dans un hameau de Cuet, chapelle vicariale de Montrevel, un jeune enfant de 9 ans et demi faisant paître le troupeau de son père. C'était Pierre Louis Marie Chanel, le futur apôtre de Futuma, le premier martyr de l'Océanie, déclaré Bienheureux par N.S. Père le Pape Léon XIII le 17 Novembre 1889.

Le 23 Mars 1817 Pierre Chanel recevait pour la première fois son Dieu des mains de Mr Trompier. Le 30 Octobre 1819 il l'envoya en classe de quatrième au petit séminaire de Méximieux.

% --- Page imprimée 179 — vue 182
Mr Trompier eut pour son élève toute la sollicitude et l'amour d'un père. Pierre Chanel conserva toujours pour celui qui avait dirigé ses premiers pas vers le sacerdoce, une reconnaissance qu'il se plaisait à manifester dans toutes ses lettres, aimant à redire que c'était à Mr Trompier, après Dieu, qu'il devait sa vocation.

XII-\underline{N... Lardet} naquit à Laye, commune d'Ouroux. Une de ses sœurs, en religion Sœur Timothée, prit une part considérable à la fondation de la communauté de St-Joseph à Monsols. Ordonné prêtre en 1790, il mourut curé de Trivy au diocèse d'Autun en 1819.

XIII-\underline{N... Ducroux} fut ordonné prêtre en 1807 et mourut curé de Balbigny (Loire). Il était né au hameau des Jambons.

XIV-\underline{N... Jambon} naquit au hameau de Braille et fut ordonné prêtre en 1820. Il est décédé en 1847 curé de St-Étienne sur Chalaronne.

XV-\underline{N... Desplaces} appartenant à la famille Desplaces du Thel. Il est mort curé de St-Appolinaire, canton de Tarare en 1845

XVI-\underline{Jean Marie Lacroix} fils de Benoît Lacroix et de Marie Janin naquit le 22 Octobre 1805 et fut baptisé le lendemain par Mr Trouilloux. La famille habitant une maison aujourd'hui disparue, sur le versant sud de la montagne de Fut-Froid. Il fut ordonné prêtre en 1830, puis fut successivement curé de Montverdun, de Saurain, de Frontenas. Il mourut aumônier de l'antiquaille en 1850

XVII-\underline{Jean Marie Bethelon} fils de Philippe Berthelon et de Jeanne Marie Lacharme naquit à Ouroux, au bourg, le 9 Décembre 1805.

Ordonné prêtre à Dijon en 1833 il fut curé à Ivry jusqu'en 1845. Il se retira à l'antiquaille à Lyon où il mourut en 1850.

XVIII-\underline{Jean Pierre Dumoulin} naquit également au bourg d'Ouroux le 14 février 1806. Son père était Benoît Dumoulin et sa mère Marie Large. Ordonné prêtre en 1830 il fut nommé vicaire à Roanne puis, curé à Chiroubles en 1839.

XIX-\underline{Jean Marie Tarlet} né au bourg d'Ouroux le 4 Décembre 1811 et baptisé le lendemain. Il eut pour parrain Jean Trompier curé de Cras. Son père, Jean Tarlet marchand à St-Pierre-le-Vieux avait épousé le 10 juillet 1810 Jeanne Marie Martin, fille de Jean Martin, cordonnier à St-Pierre et de Jeanne Marie Trompier marchande à Ouroux.

% --- Page imprimée 180 — vue 183
Mr Tarlet fut ordonné prêtre vers 1836 et nommé curé de Veziat, dans le diocèse de Belley. Après cinq ans de ministère, il entra au séminaire de Méximieux comme économe, fut nommé supérieur de la maison. Atteint de surdité il dut quitter cette charge ; il est mort avec le titre de directeur.

XX-\underline{Jean Pierre Mahuet} fils de Pierre Mahuet et de Claudine Ruet est né le 5 Mai 1824 et a été ordonné prêtre en 1847. Nommé vicaire à St-Igny-de-Vers ; il devint aumônier de l'école normale de Villefranche en 1851. Il fut nommé ensuite curé de la Bénissons-Lieu (Loire)

XXI-\underline{Jean Marie Martin} est né le 6 Avril 1829. Fils d'André Martin boucher à Ouroux et de Philiberte Lacharme, il fut ordonné prêtre en 1852. Nommé vicaire à St-Genest-Malifaux (Loire) puis curé du Bessat et enfin de St-Cyprien.

XXII-\underline{Marius Canard} fils de Jean Canard, natif du Paquier (St-Christophe) et de Jeanne Marie Aufert (fille de Jean Aufert, régisseur au château de Lacarelle et maire d'Ouroux)

Né à Lyon en 1853 il entra dans la société des prêtres de St-Irénée, fut ordonné prêtre le 26 Mai 1877, nommé professeur à l'Institution des Chartreux, vicaire de St-Bruno en 1879, poste qu'il occupa jusqu'en 1886 avant d'être missionnaire dans la même société.

XXIII-\underline{Jean Baptiste Gonnet} fils de Jean Marie Gonnet et de Madeleine Besson, né en 1881, ordonné prêtre en 1905, vicaire à St Nizier de Fornas, puis à Usson et curé à St-Georges de Renèmms et enfin de Marcy-sur-Anse.

XXIV-\underline{Gaspard Desperrier} né le 5 Février 1883, baptisé le 12, fils de Benoît Desperrier et de Jeanne Lapalud. Il a été ordonné prêtre en 1908. Depuis 1940 il est curé de Belleroche dans la Loire.

----------------

Depuis une cinquantaine d'années la Paroisse d'Ouroux n'a donné aucun prêtre à l'Église. D'où vient cette pénurie de vocations ? Les familles sont-elles restées profondément chrétiennes, assez généreuses pour donner des enfants au service de Dieu ?
